\subsection{(a) Formulation, properties and hourly decomposition}

In every hour, the consumer decides how much to consume, how much PV to use, and how much
to buy from or sell to the grid, so as to maximize its daily net utility (utility from
consuming, minus the cost of buying and producing energy, plus the revenue from selling).

\subsubsection{Sets and indices}
\begin{itemize}
    \item $t \in \T = \{0, \dots, 23\}$: the 24 hours of the day.
\end{itemize}

\subsubsection{Input data}
\begin{itemize}
    \item $p_t$: hourly electricity price in hour $t$ (DKK/kWh).
    \item $\tauimp$, $\tauexp$: import and export grid tariffs (DKK/kWh), with the
    assumption
    \begin{equation}
        \tauexp < p_t, \qquad \forall t \in \T,
        \label{eq:tariff_assumption}
    \end{equation}
    i.e.\ the export tariff is always smaller than the hourly electricity price.
    \item $\uL$: marginal consumption utility (DKK/kWh) -- a constant, the same in every
    hour.
    \item $\cpv$: PV marginal production cost (DKK/kWh) -- a constant, the same in every
    hour.
    \item $\PVmax$: maximum available PV production in hour $t$ (kWh/h).
    \item $\Lmin$, $\Lmax$: minimum and maximum hourly load consumption (kWh/h), with
    $0 \leq \Lmin \leq \Lmax$.
\end{itemize}
The consumer does not face the market price directly but the \emph{effective} import and
export prices used in the objective \eqref{eq:daily_net_utility}, which follow from the
market price and the tariffs: the import tariff is added to the market price (what the
consumer pays), and the export tariff is subtracted from it (what the consumer is paid).
\begin{align}
    \pimp &= p_t + \tauimp, & \forall t \in \T, \label{eq:effective_import_price} \\
    \pexp &= p_t - \tauexp, & \forall t \in \T. \label{eq:effective_export_price}
\end{align}

\subsubsection{Decision variables}
\begin{itemize}
    \item $\qimp \geq 0$: electricity imported from the grid in hour $t$ (kWh/h).
    \item $\qexp \geq 0$: electricity exported to the grid in hour $t$ (kWh/h).
    \item $\load$: load consumed in hour $t$ (kWh/h); its bounds are given by constraint
    \eqref{eq:c_load_bounds} below (non-negative since $\Lmin \geq 0$: a load can only
    consume, not generate, electricity).
    \item $\pvprod$: PV production used in hour $t$ (kWh/h).
\end{itemize}

\subsubsection{Objective function}
%\todo{Write the daily net utility (surplus) to maximize, as a sum over the 24 hours ($t \in
%\T$), in terms of the variables and data defined above.}
\begin{equation}
    \max_{\{\load, \qimp, \qexp, \pvprod\}_{t \in \T}} \quad
    U = \sum_{t \in \T} \Big( \uL \, \load - \pimp \, \qimp + \pexp \, \qexp - \cpv \, \pvprod \Big)
    \label{eq:daily_net_utility}
\end{equation}

\subsubsection{Constraints}

\paragraph{Power balance.}
Whatever is consumed has to come from PV production or from imports, minus whatever is
exported:
\begin{equation}
    \load = \pvprod + \qimp - \qexp, \qquad \forall t \in \T. \label{eq:c_balance}
\end{equation}

\paragraph{Load bounds.}
\begin{equation}
    \Lmin \leq \load \leq \Lmax, \qquad \forall t \in \T. \label{eq:c_load_bounds}
\end{equation}

\paragraph{PV limits.}
\begin{equation}
    0 \leq \pvprod \leq \PVmax, \qquad \forall t \in \T. \label{eq:c_pv_limit}
\end{equation}

\paragraph{Non-negativity of imports and exports.}
\begin{align}
    \qimp &\geq 0, & \forall t \in \T, \label{eq:c_import_nonneg} \\
    \qexp &\geq 0, & \forall t \in \T. \label{eq:c_export_nonneg}
\end{align}

\subsubsection{Properties of the problem}

\paragraph{Problem class and convexity.}
The objective \eqref{eq:daily_net_utility} and all the constraints
\eqref{eq:c_balance}--\eqref{eq:c_export_nonneg} are linear in the decision variables, so the
problem is a Linear Program (LP). An LP is a convex problem: its feasible set is a
polyhedron (an intersection of half-spaces and hyperplanes), and its objective is linear,
hence both concave and convex.

\paragraph{Feasibility.}
The problem is feasible as long as $\Lmin \leq \Lmax$ and $\PVmax \geq 0$ for every
$t \in \T$. Import and export are unrestricted, so nothing else can make the problem
infeasible: for any load in $[\Lmin, \Lmax]$ the balance \eqref{eq:c_balance} can be
satisfied through the grid.

\paragraph{Boundedness.}
The load and the PV production are bounded, so the objective could only be unbounded through
$\qimp$ and $\qexp$: importing and exporting large, equal amounts of energy at the same time
leaves the balance \eqref{eq:c_balance} unchanged, and would make the objective arbitrarily
large if it were profitable. This is not the case, since the effective prices satisfy
\begin{equation}
    \pexp = p_t - \tauexp < p_t + \tauimp = \pimp, \qquad \forall t \in \T,
\end{equation}
because both tariffs are positive: importing always costs more than exporting pays, so
simultaneous import and export always loses money. Hence the objective is bounded.

\paragraph{Size.}
%\todo{Number of variables and number of constraints, as a function of the 24-hour horizon.}
Each hour has four decision variables ($\load$, $\qimp$, $\qexp$, $\pvprod$), so the daily
problem has $4 \times 24 = 96$ variables. Each hour also has seven constraints: the power
balance \eqref{eq:c_balance} (an equality), the two load bounds \eqref{eq:c_load_bounds},
the non-negativity of import and export \eqref{eq:c_import_nonneg}--\eqref{eq:c_export_nonneg},
and the two PV bounds \eqref{eq:c_pv_limit}. This gives $7 \times 24 = 168$ constraints in
total, of which 24 are equalities and 144 are inequalities. If the simple bounds on the
variables are not counted as constraints, only the 24 power-balance equalities remain.

\subsubsection{Hourly decomposition}

\paragraph{Decomposition.}
%\todo{Show that the daily problem decomposes into 24 independent hourly problems: identify
%every term of the objective and every constraint, and check whether any of them links two
%different hours together.}
The daily objective \eqref{eq:daily_net_utility} is a sum over hours, and every term of
that sum only involves the variables of its own hour $t$. Likewise, every constraint
\eqref{eq:c_balance}--\eqref{eq:c_export_nonneg} only involves variables of a single hour: no
constraint links hour $t$ with any other hour. There is therefore no dependence between
different hours, and the daily problem splits into 24 independent hourly problems.

\paragraph{The hourly problem.}
%\todo{Write out the single-hour problem explicitly (its own objective and constraints, for one generic hour $t$).}
For a generic hour $t \in \T$, the problem has the same structure as the daily one without
the summation, with the data specific to that hour ($p_t$, and hence $\pimp$ and $\pexp$,
and $\PVmax$):
\begin{align}
    \max_{\load,\, \qimp,\, \qexp,\, \pvprod} \quad
    & U_t = \uL \, \load - \pimp \, \qimp + \pexp \, \qexp - \cpv \, \pvprod \\
    \text{s.t.} \quad
    & \load = \pvprod + \qimp - \qexp, \\
    & \Lmin \leq \load \leq \Lmax, \\
    & \qimp \geq 0, \quad \qexp \geq 0, \\
    & 0 \leq \pvprod \leq \PVmax.
\end{align}

\paragraph{Computational meaning.}
%\todo{What does this decomposition change about how the problem can be solved (size of the
%problem actually being solved, parallelisability, ...)?}
Since the hours are independent, the 24 hourly problems can be solved separately, in any
order, and even in parallel. Each of them is also much smaller than the daily problem: 4
variables and 7 constraints instead of 96 and 168.

This can be seen in the compact matrix form
\begin{equation}
    \min_{\mathbf{x} \in \mathbb{R}^N} \; z = \mathbf{c}^\top \mathbf{x}
    \quad \text{s.t.} \quad A\mathbf{x} \geq \mathbf{b}, \quad \mathbf{x} \geq \mathbf{0},
\end{equation}
where $\mathbf{x}, \mathbf{c} \in \mathbb{R}^N$ are the vectors of decision variables and
costs, $\mathbf{b} \in \mathbb{R}^M$ is the right-hand-side vector, and
$A \in \mathbb{R}^{M \times N}$ is the left-hand-side constraint matrix, with one row per
constraint and one column per decision variable. Ordering the variables by hour,
$\mathbf{x} = (\mathbf{x}_0, \dots, \mathbf{x}_{23})$, the absence of constraints between
different hours makes $A$ block diagonal, with one block $A_t$ per hour:
\begin{equation}
    A = \begin{pmatrix}
        A_0 & & \\
        & \ddots & \\
        & & A_{23}
    \end{pmatrix},
    \qquad
    \mathbf{c} = \begin{pmatrix} \mathbf{c}_0 \\ \vdots \\ \mathbf{c}_{23} \end{pmatrix},
    \qquad
    \mathbf{b} = \begin{pmatrix} \mathbf{b}_0 \\ \vdots \\ \mathbf{b}_{23} \end{pmatrix}.
\end{equation}
The problem therefore splits into 24 independent problems
$\min \, \mathbf{c}_t^\top \mathbf{x}_t$ s.t. $A_t \mathbf{x}_t \geq \mathbf{b}_t$,
$\mathbf{x}_t \geq \mathbf{0}$, one per hour.
\todo{Write explicitly $\mathbf{x}_t$, $\mathbf{c}_t$, $A_t$ and $\mathbf{b}_t$ for the
hourly problem above: this requires putting it in the form of the matrix formulation
(minimisation, constraints as $\geq$, variables $\geq 0$), so decide how you handle the
maximisation and the variable bounds.}

\paragraph{Economic meaning.}
%\todo{What does this decomposition say about the structure of the consumer's
%decision-making across the day?}

The load is flexible, but it is not subject to reaching a total consumption over the day:
the consumer only chooses how much to consume in each hour, between $\Lmin$ and $\Lmax$. For
example, consuming less in hour $t$ (say, not running the washing machine then) does not
oblige the consumer to consume more in another hour, and the value $\uL$ of each kWh does
not depend on when it is consumed. Therefore, to decide in hour $t$ only the data of that
hour are needed: not just the price, but also the tariffs and the PV available in that
hour, and the consumer does not need to plan ahead across the day. The flexibility of the
load does not break the independence between hours; what would break it is a minimum daily
energy requirement or a battery.

% Question 1.(b) -- Lagrange function, dual problem and KKT conditions of the hourly problem
%
% SCAFFOLD ONLY: structure and \todo prompts, no derivation. Fill each \todo with your own work.
% Task text (course_materials/Assignment_1_Instructions.html, Question 1.(b)):
%   "Formulate the Lagrange function, the dual problem, and the KKT conditions of the hourly
%    problem. Does strong duality hold? Are the KKT conditions necessary and/or sufficient
%    optimality conditions? Justify your answer from the properties stated in (a)."
% Grading (Table 2, row 1.(b)): mathematical model + writing; no code.
% Source material: Lecture 3 notes -- sections 3 (Lagrangian), 4 (Lagrange duality),
%   5 (LP duality), 6 (KKT).
%
% Called from the end of sources/1_.tex with % Question 1.(b) -- Lagrange function, dual problem and KKT conditions of the hourly problem
%
% SCAFFOLD ONLY: structure and \todo prompts, no derivation. Fill each \todo with your own work.
% Task text (course_materials/Assignment_1_Instructions.html, Question 1.(b)):
%   "Formulate the Lagrange function, the dual problem, and the KKT conditions of the hourly
%    problem. Does strong duality hold? Are the KKT conditions necessary and/or sufficient
%    optimality conditions? Justify your answer from the properties stated in (a)."
% Grading (Table 2, row 1.(b)): mathematical model + writing; no code.
% Source material: Lecture 3 notes -- sections 3 (Lagrangian), 4 (Lagrange duality),
%   5 (LP duality), 6 (KKT).
%
% Called from the end of sources/1_.tex with % Question 1.(b) -- Lagrange function, dual problem and KKT conditions of the hourly problem
%
% SCAFFOLD ONLY: structure and \todo prompts, no derivation. Fill each \todo with your own work.
% Task text (course_materials/Assignment_1_Instructions.html, Question 1.(b)):
%   "Formulate the Lagrange function, the dual problem, and the KKT conditions of the hourly
%    problem. Does strong duality hold? Are the KKT conditions necessary and/or sufficient
%    optimality conditions? Justify your answer from the properties stated in (a)."
% Grading (Table 2, row 1.(b)): mathematical model + writing; no code.
% Source material: Lecture 3 notes -- sections 3 (Lagrangian), 4 (Lagrange duality),
%   5 (LP duality), 6 (KKT).
%
% Called from the end of sources/1_.tex with \input{sources/1b_duality}.
% Whatever multipliers you introduce: add them to sources/9_nomenclat.tex (with units) and, if
% you create macros for them, put the \newcommand next to the others in main.tex.

\subsection{(b) Lagrangian, dual problem and KKT conditions}

\subsubsection{Setting up}
\paragraph{Hourly problem in the form used here.}
\todo{Restate the hourly problem of (a) for a generic hour $t$ (one line, refer to the
equations of (a)). Decide and state your conventions: do you keep it as a maximisation or
turn it into a minimisation as in the lecture notes? How is each constraint written
($\leq 0$, $\geq 0$, $= 0$)? The signs of all multipliers below depend on these choices, so
fix them once here.}

\paragraph{Multipliers.}
\todo{Associate one dual variable to each constraint: give a table or list (name, the
constraint it belongs to, its sign restriction, its unit). Decide which constraints get a
multiplier and which ones stay as variable bounds ($x \geq 0$) -- say which and why.}

\subsubsection{Lagrange function}
\todo{Write the Lagrangian of the hourly problem with your multipliers.}

\subsubsection{Dual problem}
\paragraph{Dual function.}
\todo{Write the dual function (minimise/maximise the Lagrangian over the primal variables).
State the condition on the multipliers for which it is finite, and where it comes from.}

\paragraph{Dual problem.}
\todo{Write the dual problem: objective, constraints on the multipliers. Which problem class
is it?}

\subsubsection{KKT conditions}
\todo{Write the four groups of conditions for your problem, one line each per variable or
constraint: (i) stationarity, (ii) primal feasibility, (iii) dual feasibility, (iv)
complementary slackness.}

\subsubsection{Strong duality}
\todo{Does strong duality hold? Justify from the properties of the problem stated in (a)
(problem class, feasibility, boundedness), and state what it implies for the primal and dual
optimal values.}

\subsubsection{Necessity and sufficiency of the KKT conditions}
\todo{Are the KKT conditions necessary, sufficient, both, or neither for optimality of this
problem? Justify from the properties stated in (a), separately for necessity and for
sufficiency.}
.
% Whatever multipliers you introduce: add them to sources/9_nomenclat.tex (with units) and, if
% you create macros for them, put the \newcommand next to the others in main.tex.

\subsection{(b) Lagrangian, dual problem and KKT conditions}

\subsubsection{Setting up}
\paragraph{Hourly problem in the form used here.}
\todo{Restate the hourly problem of (a) for a generic hour $t$ (one line, refer to the
equations of (a)). Decide and state your conventions: do you keep it as a maximisation or
turn it into a minimisation as in the lecture notes? How is each constraint written
($\leq 0$, $\geq 0$, $= 0$)? The signs of all multipliers below depend on these choices, so
fix them once here.}

\paragraph{Multipliers.}
\todo{Associate one dual variable to each constraint: give a table or list (name, the
constraint it belongs to, its sign restriction, its unit). Decide which constraints get a
multiplier and which ones stay as variable bounds ($x \geq 0$) -- say which and why.}

\subsubsection{Lagrange function}
\todo{Write the Lagrangian of the hourly problem with your multipliers.}

\subsubsection{Dual problem}
\paragraph{Dual function.}
\todo{Write the dual function (minimise/maximise the Lagrangian over the primal variables).
State the condition on the multipliers for which it is finite, and where it comes from.}

\paragraph{Dual problem.}
\todo{Write the dual problem: objective, constraints on the multipliers. Which problem class
is it?}

\subsubsection{KKT conditions}
\todo{Write the four groups of conditions for your problem, one line each per variable or
constraint: (i) stationarity, (ii) primal feasibility, (iii) dual feasibility, (iv)
complementary slackness.}

\subsubsection{Strong duality}
\todo{Does strong duality hold? Justify from the properties of the problem stated in (a)
(problem class, feasibility, boundedness), and state what it implies for the primal and dual
optimal values.}

\subsubsection{Necessity and sufficiency of the KKT conditions}
\todo{Are the KKT conditions necessary, sufficient, both, or neither for optimality of this
problem? Justify from the properties stated in (a), separately for necessity and for
sufficiency.}
.
% Whatever multipliers you introduce: add them to sources/9_nomenclat.tex (with units) and, if
% you create macros for them, put the \newcommand next to the others in main.tex.

\subsection{(b) Lagrangian, dual problem and KKT conditions}

\subsubsection{Setting up}
\paragraph{Hourly problem in the form used here.}
\todo{Restate the hourly problem of (a) for a generic hour $t$ (one line, refer to the
equations of (a)). Decide and state your conventions: do you keep it as a maximisation or
turn it into a minimisation as in the lecture notes? How is each constraint written
($\leq 0$, $\geq 0$, $= 0$)? The signs of all multipliers below depend on these choices, so
fix them once here.}

\paragraph{Multipliers.}
\todo{Associate one dual variable to each constraint: give a table or list (name, the
constraint it belongs to, its sign restriction, its unit). Decide which constraints get a
multiplier and which ones stay as variable bounds ($x \geq 0$) -- say which and why.}

\subsubsection{Lagrange function}
\todo{Write the Lagrangian of the hourly problem with your multipliers.}

\subsubsection{Dual problem}
\paragraph{Dual function.}
\todo{Write the dual function (minimise/maximise the Lagrangian over the primal variables).
State the condition on the multipliers for which it is finite, and where it comes from.}

\paragraph{Dual problem.}
\todo{Write the dual problem: objective, constraints on the multipliers. Which problem class
is it?}

\subsubsection{KKT conditions}
\todo{Write the four groups of conditions for your problem, one line each per variable or
constraint: (i) stationarity, (ii) primal feasibility, (iii) dual feasibility, (iv)
complementary slackness.}

\subsubsection{Strong duality}
\todo{Does strong duality hold? Justify from the properties of the problem stated in (a)
(problem class, feasibility, boundedness), and state what it implies for the primal and dual
optimal values.}

\subsubsection{Necessity and sufficiency of the KKT conditions}
\todo{Are the KKT conditions necessary, sufficient, both, or neither for optimality of this
problem? Justify from the properties stated in (a), separately for necessity and for
sufficiency.}
